%\myChapterStar{Titre}{Titre court}{Ajouter a la table des matieres? (false|true|chapter|section|subsection|subsubsection -chapter par defaut-)}
\myChapterStar{Conclusion g\'en\'erale}{}{true}
\label{chap:conclusion}
%\myMinitoc{Profondeur de la minitoc (section|subsection|subsubsection)}{Titre de la minitoc}
\myMiniToc{section}{Sommaire}

Ce m\'emoire a abord\'e la pr\'ediction du risque de d\'epassement r\'eglementaire PFAS dans les eaux souterraines en mode pr\'edictif strict, \`a partir de variables contextuelles g\'eospatiales seules et sans recours \`a des concentrations pr\'ealablement mesur\'ees. Nous rappelons ci-apr\`es la probl\'ematique et les choix m\'ethodologiques qui ont structur\'e ce travail, avant d'en analyser les r\'esultats de mani\`ere critique et de tracer les perspectives qu'ils ouvrent.

%\mySectionStar{Titre}{Titre court}{Ajouter a la table des matieres? (false|true|chapter|section|subsection|subsubsection -section par defaut-)}
\mySectionStar{Probl\'ematique \'etudi\'ee et choix m\'ethodologiques}{}{true}
\label{sec:concl-problematique}

La surveillance de la qualit\'e des eaux souterraines face \`a la contamination PFAS pose un d\'efi op\'erationnel concret~: l'analyse chimique par chromatographie liquide coupl\'ee \`a la spectrom\'etrie de masse co\^ute plusieurs centaines d'euros par \'echantillon, et les limites fix\'ees par l'EPA en 2024 dans le cadre de la r\'eglementation NPDWR descendent jusqu'\`a 4~ng/L pour le PFOA et le PFOS~\cite{epa2023npdwr}. Pr\'edire o\`u ces seuils sont susceptibles d'\^etre d\'epass\'es, sans concentration PFAS pr\'ealablement mesur\'ee, constitue donc un besoin r\'eel pour la priorisation des campagnes d'\'echantillonnage. Ce travail s'est attaqu\'e \`a ce probl\`eme sur un corpus de 46\,338 pr\'el\`evements d'eaux souterraines californiennes, constitu\'e en s'inspirant de la m\'ethodologie de collecte de Dong et al.~\cite{dong2024pfas}. Deux verrous ont structur\'e la d\'emarche.

Le premier est d'ordre logiciel. Toute approche de pr\'ediction g\'eospatiale suppose que les observations soient enrichies par des variables de contexte, telles que la proximit\'e des sources de contamination, les propri\'et\'es des sols et les variables hydrog\'eologiques, car ces caract\'eristiques sont absentes des jeux de donn\'ees PFAS bruts~\cite{dong2024pfas,george2021pfas}. Or cet enrichissement est g\'en\'eralement assur\'e par des scripts monolithiques, \'etroitement li\'es \`a un jeu de donn\'ees particulier et non r\'eutilisables. Pour lever ce verrou, nous avons con\c{c}u \texttt{geoenrich}, un moteur g\'en\'erique organis\'e en trois couches~: une couche \emph{contrat} agnostique, une couche \emph{op\'erateurs} rassemblant des primitives de jointure param\'etrables, et une couche \emph{configuration} qui concentre les choix propres \`a chaque r\'egion dans un fichier d\'eclaratif externe. Sa correction a \'et\'e \'etablie par un test de parit\'e avec le pipeline d'origine~: 77 colonnes strictement identiques, 17 mieux renseign\'ees et aucune r\'egression sur 94 colonnes produites.

Le second verrou est scientifique et tient \`a l'\'evaluation des mod\`eles. La contamination PFAS \'etant fortement auto-corr\'el\'ee dans l'espace~\cite{guelfo2021pfas}, un d\'ecoupage al\'eatoire train/test laisse des forages spatialement voisins de part et d'autre de la fronti\`ere, permettant \`a un mod\`ele de m\'emoriser la g\'eographie plut\^ot que d'apprendre des m\'ecanismes transf\'erables. Pour y r\'epondre, toute variable de localisation a \'et\'e retir\'ee des entr\'ees des mod\`eles~; les coordonn\'ees ne subsistent que dans la topologie du graphe relationnel construit pour le HGT. Cinq mod\`eles ont \'et\'e entra\^in\'es et compar\'es \`a protocole identique sur ce jeu enrichi~: deux mod\`eles tabulaires de r\'ef\'erence (for\^et al\'eatoire et XGBoost), un transformeur de graphe h\'et\'erog\`ene (HGT), et deux mod\`eles hybrides (Embedding Fusion et Stacking). Chacun a \'et\'e accompagn\'e d'analyses d'interpr\'etabilit\'e globale et locale par valeurs de Shapley~\cite{lundberg2017shap}.

\mySectionStar{Analyse critique des r\'esultats obtenus}{}{true}
\label{sec:concl-analyse}

Les r\'esultats obtenus sur le jeu de test, en mode pr\'edictif strict, montrent que les mod\`eles tabulaires dominent. La for\^et al\'eatoire et le stacking atteignent tous deux un score F1 de 0,925, avec des aires sous la courbe ROC de 0,979 et 0,977~; XGBoost suit de tr\`es pr\`es (F1 0,923, AUC 0,977). Le HGT relationnel, qui ne re\c{c}oit la g\'eographie que par la structure du graphe, reste nettement en retrait (F1 0,825, AUC 0,906), et la fusion par embeddings se situe entre les deux familles (F1 0,887, AUC 0,960). L'enseignement principal de cette comparaison est que les variables de contexte produites par \texttt{geoenrich}, m\^eme priv\'ees de coordonn\'ees, portent l'essentiel du signal pr\'edictif~: des mod\`eles tabulaires classiques atteignent une discrimination \'elev\'ee sans acc\`es \`a la localisation. Le graphe relationnel capte un signal r\'eel, mais il reste en retrait des mod\`eles tabulaires. Le stacking, arbitr\'e par un m\'eta-classifieur XGBoost, n'exploite pas la contribution du HGT et n'am\'eliore pas la for\^et al\'eatoire~: ce r\'esultat indique qu'un enrichissement de contexte soign\'e peut rendre superflue une mod\'elisation relationnelle plus lourde.

Ces conclusions doivent \^etre lues avec deux r\'eserves. La premi\`ere concerne l'\'evaluation~: les performances rapport\'ees reposent sur un d\'ecoupage strat\'ifi\'e al\'eatoire, qui ne teste pas la g\'en\'eralisation \`a des zones enti\`erement non vues. La contamination PFAS \'etant auto-corr\'el\'ee dans l'espace, une partie des scores pourrait tenir \`a une proximit\'e r\'esiduelle entre forages d'entra\^inement et de test plut\^ot qu'\`a des m\'ecanismes transf\'erables. Les chiffres obtenus constituent donc une borne sup\'erieure optimiste~; une validation crois\'ee spatiale par blocs permettrait de mesurer cet \'ecart. La seconde r\'eserve concerne \texttt{geoenrich}~: son ind\'ependance vis-\`a-vis du cas californien est \'etablie par inspection du code et par conception, mais la preuve empirique d\'ecisive reste \`a apporter. Ces deux limites ne remettent pas en cause la comparaison entre mod\`eles, qui s'effectue \`a protocole identique et demeure valide en termes relatifs.

L'analyse de l'interpr\'etabilit\'e globale par valeurs de Shapley confirme que les variables les plus influentes sont des grandeurs de contexte hydrog\'eologique et de proximit\'e aux sources (cat\'egorie du forage, profondeur de la nappe, densit\'e de sites de contamination dans le voisinage, direction d'\'ecoulement), et non des coordonn\'ees. Ce r\'esultat est coh\'erent avec le dispositif de retrait de la localisation et l\'egitime r\'etrospectivement la d\'ecision de ne pas fournir les coordonn\'ees comme variables d'entr\'ee.

\mySectionStar{Quelques perspectives}{}{true}
\label{sec:concl-perspectives}

Ce travail ouvre plusieurs directions de recherche et de d\'eveloppement que nous estimons prioritaires.

\begin{itemize}

    \item \textbf{Validation crois\'ee spatiale par blocs.} Cette validation n'a pas \'et\'e conduite dans le pr\'esent travail. Elle constitue la priorit\'e imm\'ediate~: elle permettrait de distinguer la part du signal pr\'edictif qui reste stable hors des zones d'entra\^inement de celle qui repose sur l'auto-corr\'elation spatiale r\'esiduelle, et de disposer ainsi d'une mesure honn\^ete de la transf\'erabilit\'e g\'eographique des mod\`eles.

    \item \textbf{D\'emonstration empirique de la g\'en\'ericit\'e de \texttt{geoenrich}.} L'ind\'ependance du moteur vis-\`a-vis du cas PFAS californien est \'etablie par construction. Convertir cette propri\'et\'e en preuve d\'emontr\'ee suppose de soumettre \`a l'outil un jeu d'observations portant sur un autre contaminant (nitrates, m\'etaux lourds, pesticides) ou une autre r\'egion, avec des couches de r\'ef\'erence adapt\'ees, et de v\'erifier que l'enrichissement produit est pertinent sans modifier une ligne du code du moteur. Ce test de g\'en\'eralit\'e consoliderait la contribution logicielle de ce m\'emoire.

    \item \textbf{Int\'egration de l'interpr\'etabilit\'e dans des outils d'aide \`a la d\'ecision.} Les explications locales par valeurs de Shapley~\cite{lundberg2017shap} permettent, pour chaque forage, d'identifier les facteurs qui ont orient\'e la pr\'ediction vers le d\'epassement ou vers la conformit\'e. Leur int\'egration dans un tableau de bord destin\'e aux gestionnaires de r\'eseaux d'eau rendrait les pr\'edictions lisibles et actionnables, en particulier pour la hi\'erarchisation des campagnes de contr\^ole analytique.

    \item \textbf{Application au Cameroun et au contexte africain.} La probl\'ematique pos\'ee par ce travail d\'epasse le cadre californien. Dans de nombreux pays africains, dont le Cameroun, les r\'eseaux de surveillance de la qualit\'e des eaux souterraines restent peu denses, et le recours syst\'ematique \`a l'analyse chimique des PFAS demeure hors de port\'ee budg\'etaire pour la plupart des gestionnaires. C'est pr\'ecis\'ement dans ces contextes de raret\'e des donn\'ees que le mode pr\'edictif strict pr\'esente le plus de valeur~: il ne requiert aucune concentration PFAS pr\'ealable et s'appuie uniquement sur le contexte g\'eospatial (proximit\'e des sites militaires et industriels ayant utilis\'e des mousses anti-incendie AFFF, propri\'et\'es des sols, variables hydrog\'eologiques). \texttt{geoenrich} a \'et\'e con\c{c}u pour \^etre r\'egion-agnostique~: assembler les couches de r\'ef\'erence pertinentes \`a partir de bases mondiales (ISRIC SoilGrids pour les sols, bases de donn\'ees hydrog\'eologiques continentales, registres de sites industriels nationaux) suffirait, en principe, \`a l'instancier pour un contexte camerounais ou africain. Une telle extension supposerait de collecter un \'echantillon de mesures de calibration sur le terrain, de constituer les couches de r\'ef\'erence correspondantes et de valider le pipeline dans ce nouvel environnement. Les r\'esultats obtenus sur le dataset californien constituent un point de d\'epart encourageant pour envisager cette transposition.

\end{itemize}


\myCleanStarChapterEnd
